\section{The Analytical Framework: Formalizing the
	Transformation Elasticity}
\label{sec:3:analytical_framework}

Distributive pressure conditions the choice of technique, and the choice
of technique determines how much capacity a unit of accumulation
produces. Section~2 established that claim historically and
institutionally. It carries no elasticity until the cost-minimization
problem behind it is written down, because only then does $\theta$ acquire
a determinate value rather than a direction of movement.
\citet{CajasGuajirro2024} first posed that problem as a choice over
mechanization growth, and I extend it in three steps.

The steps are nested rather than alternative. A firm that treats capital
as homogeneous faces a single margin, and the elasticity that emerges
from its cost minimization is a scalar. Allow the capital stock to
separate into nonresidential structures and machinery, and a second
margin opens: the mechanization gap $\kappa_t$, whose marginal capacity
payoff depends on distribution and on relative asset prices at once.
Impose a foreign-exchange bottleneck on the machinery half of that gap,
and the interior solution can fail entirely, truncating $\theta$ at a
corner. Setting the heterogeneity parameter to zero recovers the first
model from the second; relaxing the foreign-exchange constraint recovers
the center from the periphery. One structure, evaluated under three
constraint sets.

A limitation is worth stating before the derivations rather than after
them. The mechanism is multi-sectoral in principle---disproportionality
is a relation between departments of production---while the empirical
implementation bounds it to a single-sector aggregate. That bound masks
internal sectoral divergence. It preserves the transformation logic,
which is what the estimation requires, but it means $\theta$ is recovered
as an economy-wide average of a relation that operates unevenly across
branches.

\subsection{The Baseline Model of Cost-Minimizing
	Mechanization (Homogeneous Capital Stock)}

Actual aggregate output is represented by a Leontief
production function with fixed technical coefficients:
\begin{equation}
	Y = \min\left( A L,\; \mu B K \right)
	\label{eq:31:leontief}
\end{equation}
where $A = Y/L$ is actual labor productivity, $\mu = Y/Y^p$
is capacity utilization, and $B = Y^p/K$ is maximum capital
productivity. This production structure governs the rate of
profit through the fundamental heterodox accounting identity:
\begin{equation}
	r = \mu B (1 - \omega)
	\label{eq:31:profit_identity}
\end{equation}
where $\omega$ is the wage share. Standard macroeconomic
models treat $B$ as an exogenous technological given and $\mu$
as a behavioral constant smoothly targeted by a unified firm.
I reject both assumptions. The components of
\eqref{eq:31:profit_identity} are endogenous outcomes of
class struggle and decentralized market anarchy, as
established in Section~2. This subsection formalizes how
distributive pressure forces cost-minimizing firms to alter
their choice of technique, endogenously determining $B$.

\subsubsection{Tactical Choice of Technique and Induced
	Innovation}

Capitalist firms choice of technique is given by aiming a mechanization growth rate $q$ under competitive cutting cost-pressures, looking at the net productivity surplus of technical change as their objective function.  Let $a$ denote the growth rate of
labor productivity. The effective distributive cost of
deploying mechanization growth is $q(1-\omega)$. Hence, the choice of technique of capitalist mechanization can be represented by: 
\begin{equation}
	\min_{q}\; c(q) = q(1-\omega) - a 
	\label{eq:31:objective}
\end{equation}
This objective captures the induced-innovation logic: firms
accelerate mechanization to expand labor productivity growth,
but the distributive cost of capital intensity holds them
back. The cost-minimizing objective function of capitalist becomes constrained by a mechanization function that can be represented as:

\begin{equation}
	a = \Phi(q) = \lambda_0 + \lambda_1 q + \lambda_2 q^2
	\label{eq:31:frontier}
\end{equation}

\textbf{Assumption 1.} \textit{Concavity of the
	mechanization frontier.} The frontier parameters satisfy
$\lambda_0 > 0$, $\lambda_1 > 0$, and $\lambda_2 < 0$. The
base productivity growth $\lambda_0$ sets the floor without
advanced machinery. The linear coefficient $\lambda_1$
captures the direct productivity payoff of mechanization. The
strict concavity $\lambda_2 < 0$ ensures diminishing marginal
returns to mechanization.

Substituting \eqref{eq:31:frontier} into
\eqref{eq:31:objective} and solving the first-order condition
yields the optimal mechanization growth:
\begin{align}
	c(q) &= \lambda_0 + \left[\lambda_1 - (1-\omega)\right] q
	+ \lambda_2 q^2 \label{eq:31:unconstrained} \\
	\frac{\partial c}{\partial q} &= \lambda_1 - (1-\omega) + 2\lambda_2 q^*
	= 0 \label{eq:31:foc} \\
	q^*(\omega) &= \frac{(1-\omega) - \lambda_1}{2\lambda_2}
	\label{eq:31:qstar}
\end{align}
The strict concavity ($c''(q) = 2\lambda_2 < 0$) guarantees
a unique maximum.

\textbf{Proposition 1.} \textit{Induced innovation.} Under
Assumption~1, the optimal mechanization growth is strictly
increasing in the wage share:
\begin{equation}
	\frac{\partial q^*}{\partial \omega}
	= -\frac{1}{2\lambda_2} > 0
	\label{eq:31:comparative_static}
\end{equation}
A rise in the wage share forces cost-minimizing firms to
substitute machinery for labor. Rising wages drive
mechanization directly through the first-order condition,
pushing optimal capital intensity growth upward.
\textit{Proof.} Immediate from differentiating
\eqref{eq:31:qstar}. $\hfill\blacksquare$

\textbf{Remark 1.} The objective
\eqref{eq:31:objective} does not assert that the firm is a
coherent optimizer. Following \citet{Vidal2022}, the firm is
internally divided between the \textit{control imperative}
(valorization: cost reduction through mechanization) and the
\textit{cooperation imperative} (production: the need for
worker effort, tacit knowledge, and problem-solving
capacity). The formal model captures the cost-minimizing
pole---the Okishio viability logic---while the constraints on
the optimization problem capture the limits imposed by the
cooperation imperative and by external conditions. The firm
maximizes net productivity surplus---labor productivity
growth minus the distributive cost of mechanization---subject
to diminishing returns. The firm's actual mechanization
choice is a contested settlement between these imperatives,
not a coherent optimum. What the model demonstrates is the
cost-reducing direction in which distributive pressure pushes
the settlement, not the settlement itself.

\subsubsection{Endogenous Transformation Elasticity}

Substituting the optimal technique back into the
technological frontier recovers the normal growth rate of
labor productivity under smooth reproduction:
\begin{equation}
	a^N(\omega) = \lambda_0 + \lambda_1 q^*(\omega)
	+ \lambda_2 \left[ q^*(\omega) \right]^2
	\label{eq:31:aN}
\end{equation}
This optimal technique dictates the structural trajectory of
the Leontief parameter $B$. By definition, maximum capital
productivity is $B = Y^p / K$. In dynamic growth rates, the
growth of capital productivity is:
\begin{equation}
	b \equiv g_{Y^p} - g_K = a - q
	\label{eq:31:b_def}
\end{equation}
Evaluated at the smooth reproduction benchmark:
$b^N(\omega) = a^N(\omega) - q^*(\omega)$.

\textbf{Definition 1.} \textit{Transformation elasticity.}
The transformation elasticity $\theta$ is the structural
parameter that translates capital accumulation into productive
capacity growth:
\begin{equation}
	g_{Y^p} = \theta\, g_K
	\label{eq:31:theta_def}
\end{equation}
Substituting \eqref{eq:31:theta_def} into
\eqref{eq:31:b_def} reveals the link between the level
coefficient $B$ and the elasticity $\theta$:
\begin{equation}
	\theta = 1 + \frac{b}{g_K}
	\label{eq:31:theta_link}
\end{equation}

To evaluate this elasticity, I anchor the choice of technique
to a baseline rate of capital accumulation. Following the
normal reproduction path, the normal accumulation rate is
driven by the normal profit rate and the normal rate of re-capitalization  net of depreciation. I follow \citet{Foley1986} definition of the re-capitalization rate, $\chi=\frac{I}{\Pi}$, and interpret its normal conditions as continuous realization and re-investment of the surplus.\footnote{A simple accounting decomposition can demonstrate that the multiplier of profitability that drives investment, is precisely $\chi$. In a monetary economy with credit, the re-capitalization rate does not require to be bounded by the surplus $\Pi$ order of magnitude.} Hence, 

\begin{equation}
	g_K^N = \chi^N r^N - \delta
	\label{eq:31:gKN}
\end{equation}
In dynamic growth accounting, changes in the growth rate of capital represent an acceleration term $\ddot{k} \equiv d g_K / dt$. In historical time, $\ddot{k}$ is kept strictly bounded by corporate liquidity constraints, debt-servicing ceilings, and Central Bank credit cycles, preventing explosive capital growth divergence. Expressed dynamically relative to this normal path, the aggregate transformation elasticity becomes:
\begin{equation}
	\theta(\omega) = 1 + \frac{b^N(\omega)}{g_K^N}
	= 1 + \frac{a^N(\omega) - q^*(\omega)}{g_K^N}
	\label{eq:31:theta_omega}
\end{equation}

\textbf{Proposition 2.} \textit{Endogeneity of the
	transformation elasticity.} Under Assumption~1, the
transformation elasticity $\theta$ is not an exogenous
technical constant. It emerges endogenously from the
cost-minimizing choice of technique. When distributive
pressure forces firms to alter mechanization via
\eqref{eq:31:comparative_static}, it simultaneously shifts
the structural efficiency of accumulation via
\eqref{eq:31:theta_omega}. The elasticity $\theta$ is
therefore a function of the wage share $\omega$, the frontier
parameters $\lambda_0, \lambda_1, \lambda_2$, and the normal
accumulation rate $g_K^N$.

\paragraph{Econometric note on capital measurement error.}
Because the aggregate capital stock $K_t$ is constructed via Perpetual Inventory (GPIM) methods involving service-life and discard schedules, any classical right-hand-side measurement error in $K_t$ introduces standard econometric attenuation bias that pulls the estimated elasticity $\hat{\theta}$ downward toward zero. Classical measurement noise cannot account for the empirical finding that $\theta < 1$; rather, it reinforces that the true structural transformation elasticity is bounded strictly below unity even under conservative measurement assumptions.


\subsection{Heterogeneous Capital, the Mechanization Gap,
	and State-Dependent Viability}

Mainstream macroeconomics treats the capital stock as a
homogeneous aggregate, abstracting from the internal
architecture of the firm. Following the theory of organizational political economy and the internally divided firm \citep{Vidal2022}, I unbundle
aggregate productive capital into and extensive margin given by nonresidential
structures ($K^{NRC}$) and intensive one, given by the accumulation of machinery and equipment ($K^{ME}$).
These components do not unfold throughout the labor process
symmetrically. Structures form the built environments that setup the 
``envelope''of accumulation --- i.e. the factory floor and logistics shell of the productive process. Capital accumulation in Machinery and Equipment 
is the active component that restructures the production process and the extraction of 
relative surplus value \citep[Vol.~I, Ch.~15]{MarxKVI}.

\textbf{Assumption 2.} \textit{Heterogeneous capital
	decomposition.} Aggregate productive capital decomposes as
$K^{cap}_t = K^{NRC}_t + K^{ME}_t$. Nonresidential
structures evolve independently of short-run distributive
conflict ($\omega$) and its rate of accumulation of accumulation is governed by its long-run relation with the rate of accumulation in machinery and equipment under normal conditions of accumulation in the long-run 
\begin{equation}
	g_{K^{NRC}, t}^N = f(\chi^N,r^N_t,g_{K^{ME}, t}^N,\delta) 
	\label{eq:32:gKNRC}
\end{equation}
This extensive envelope $\Theta(K^{NRC})$ becomes a baseline for aggregate elasticity of transformation $\theta$ such that the rate of accumulation of capital stock in non-residential construction is a byproduct of the intensive margin of accumulation decided by a tactical technique choice made by capitalists in response to distributive pressures:

\subsubsection{The Extensive Envelope and the Intensive Margin}

Given the fixed infrastructural envelope, the
\textbf{intensive margin} of accumulation is the tactical
choice of the mechanization gap, defined in natural logs as:
\begin{equation}
	\kappa_t = k^{ME}_t - k^{NRC}_t
	\label{eq:32:kappa_def}
\end{equation}
This gap measures mechanization intensity relative to the
extensive plant scale. The capacity identity decomposes
potential output into the extensive envelope and the
intensive capacity yield $\Psi_t$:
\begin{equation}
	y^p_t = \theta_{NRC} k^{NRC}_t
	+ \Psi_t(\kappa_t, \omega_{t-1}, p^\kappa_t)
	\label{eq:32:capacity_decomp}
\end{equation}
The intensive yield $\Psi_t$ is not a fixed engineering
constant. It is the outcome of a decentralized
cost-minimization problem governed by the regime of
accumulation, as formalized below.

\subsubsection{Cost-Surplus Minimization and the Okishio
	Viability Yield}

In this case, the capitalist firm chooses a mechanization gap $\kappa_t$ under cutting-cost pressures, subject to the
mechanization frontier:
\begin{equation}
	\Phi(\kappa_t) = \lambda_0 + \lambda_1 \kappa_t
	+ \lambda_2 \kappa_t^2
	\label{eq:32:frontier_kappa}
\end{equation}
The \textit{realized} capacity yield is governed by a viability condition: a
cost-minimizing firm deploys a new technique only if it
reduces unit costs at current distributive and price
parameters \citep{Okishio1961, Basu2022}. I formalize
this as a cost-minimization as problem:
\begin{equation}
	\min_{\kappa_t}\; c(\kappa_t)
	= \kappa_t \cdot MC(\omega_{t-1}, p^\kappa_t) -\Phi(\kappa_t)
	\label{eq:32:objective_kappa}
\end{equation}

The marginal viability cost $MC(\omega_{t-1}, p^\kappa_t)$
is determined by two channels:
\begin{enumerate}
	\item \textbf{The distributive dividend (labor saving).}
	Increasing $\kappa_t$ substitutes machinery for labor. The
	value of this displacement is proportional to the lagged
	wage share ($\omega_{t-1}$), reflecting the wage
	settlement period prior to current capacity realization.
	A higher wage share \textit{reduces} the net cost of
	mechanization.
	\item \textbf{The relative price penalty (capital cost).}
	Increasing $\kappa_t$ requires accumulating machinery
	faster than structures. The cost of this composition shift
	is proportional to the log relative price of machinery
	($p^\kappa_t = \ln P^{ME}_t - \ln P^{NRC}_t$). A higher
	relative price \textit{increases} the marginal cost.
\end{enumerate}
I specify the marginal viability cost as:
\begin{equation}
	MC(\omega_{t-1}, p^\kappa_t)
	= \gamma_0 - \gamma_1 \omega_{t-1} + \gamma_2 p^\kappa_t
	\label{eq:32:MC}
\end{equation}
where $\gamma_1 > 0$ captures the wage-savings dividend and
$\gamma_2 > 0$ captures the relative price penalty.

Substituting \eqref{eq:32:frontier_kappa} and
\eqref{eq:32:MC} into \eqref{eq:32:objective_kappa} and
solving the first-order condition yields the optimal
mechanization gap:
\begin{align}
	c(\kappa_t) &= \lambda_0
	+ \left[\lambda_1 - (\gamma_0 - \gamma_1 \omega_{t-1}
	+ \gamma_2 p^\kappa_t)\right] \kappa_t
	+ \lambda_2 \kappa_t^2
	\label{eq:32:unconstrained_kappa} \\
	\frac{\partial c}{\partial \kappa_t} &= \lambda_1 - \gamma_0
	+ \gamma_1 \omega_{t-1} - \gamma_2 p^\kappa_t
	+ 2\lambda_2 \kappa_t^* = 0
	\label{eq:32:foc_kappa} \\
	\kappa_t^*(\omega_{t-1}, p^\kappa_t)
	&= \frac{(\gamma_0 - \lambda_1) - \gamma_1 \omega_{t-1}
		+ \gamma_2 p^\kappa_t}{2\lambda_2}
	\label{eq:32:kappa_star}
\end{align}

\textbf{Proposition 3.} \textit{Dual-channel induced
	innovation.} Under Assumption~1 (concavity) and
Assumption~2 (heterogeneous capital), the optimal
mechanization gap satisfies:
\begin{align}
	\frac{\partial \kappa_t^*}{\partial \omega_{t-1}}
	&= -\frac{\gamma_1}{2\lambda_2} > 0
	\label{eq:32:cs_omega} \\
	\frac{\partial \kappa_t^*}{\partial p^\kappa_t}
	&= \frac{\gamma_2}{2\lambda_2} < 0
	\label{eq:32:cs_price}
\end{align}
A rising wage share forces cost-minimizing firms to widen
the mechanization gap, substituting machinery for labor.
A rising relative price of machinery suppresses the optimal
gap, as the structural cost of the composition shift
penalizes mechanization. \textit{Proof.} Immediate from
differentiating \eqref{eq:32:kappa_star}. $\hfill\blacksquare$

\subsubsection{The Marginal Capacity Payoff and the
	Econometric Bridge}

The macroeconomic translation from the mechanization gap to
potential capacity is given by the \textit{marginal capacity payoff} ($MP_\kappa$), the derivative of the
intensive capacity yield $\Psi_t$ with respect to $\kappa_t$.
By the envelope theorem:
\begin{align}
	MP_\kappa &= \frac{\partial \Psi_t}{\partial \kappa_t}
	= \frac{\partial \Phi}{\partial \kappa_t}
	- MC(\omega_{t-1}, p^\kappa_t)
	\label{eq:32:MP_env} \\
	&= (\lambda_1 + 2\lambda_2 \kappa_t)
	- (\gamma_0 - \gamma_1 \omega_{t-1} + \gamma_2 p^\kappa_t)
	\label{eq:32:MP_expanded} \\
	&= \underbrace{(\lambda_1 - \gamma_0
		+ 2\lambda_2 \kappa_t)}_{\text{autonomous yield}}
	+ \gamma_1 \omega_{t-1} - \gamma_2 p^\kappa_t
	\label{eq:32:MP_rearranged}
\end{align}

Integrating the marginal payoff with respect to $\kappa_t$
(treating $\omega_{t-1}$ and $p^\kappa_t$ as predetermined
shifters) and defining $\theta_\kappa$,
$\theta_{\omega\kappa} \equiv \gamma_1$, and
$\theta_{\kappa p\kappa} \equiv -\gamma_2$ yields the
intensive capacity yield function:
\begin{equation}
	\Psi_t = \theta_\kappa \kappa_t
	+ \theta_{\omega\kappa} (\omega_{t-1} \kappa_t)
	+ \theta_{\kappa p\kappa} (p^\kappa_t \kappa_t)
	\label{eq:32:Psi}
\end{equation}
Substituting \eqref{eq:32:Psi} back into
\eqref{eq:32:capacity_decomp} recovers the structural
capacity equation:
\begin{equation}
	y^p_t = \theta_{NRC} k^{NRC}_t + \theta_\kappa \kappa_t
	+ \theta_{\omega\kappa} (\omega_{t-1} \kappa_t)
	+ \theta_{\kappa p\kappa} (p^\kappa_t \kappa_t)
	\label{eq:32:capacity_eq}
\end{equation}
The interaction terms are not ad-hoc econometric additions.
They are the macroeconomic footprint of the divided firm's
Okishio viability constraint. Distributive pressure and
relative price enter because they alter the marginal
capacity payoff of the mechanization gap.

\subsubsection{The Transformation Elasticity under Heterogeneous Capital Stock}

% AUTHOR NOTE (not typeset): this subsection requires editing and careful
% theoretical revision.

Specification~B recovers separate capacity elasticities for nonresidential structures ($\theta_{\text{NRC}}$) and machinery ($\theta_\kappa$), but evaluating aggregate capacity dynamics requires defining an aggregate transformation elasticity ($\theta^{\text{Aggregate}}_t$). In a single-asset model (Specification~A), $\theta_A = \partial y^p / \partial k^{\text{cap}}$ is a scalar. Under heterogeneous capital, the total differential of productive capacity \eqref{eq:32:capacity_eq} separates extensive expansion from intensive composition shifters:
\begin{align}
	dy^p_t &= \theta_{\text{NRC}}\, dk^{\text{NRC}}_t + MP_\kappa\, d\kappa_t
	\label{eq:32:dy_total} \\
	&= \theta_{\text{NRC}}\, dk^{\text{NRC}}_t
	+ \left(\theta_\kappa + \theta_{\omega\kappa}\, \omega_{t-1}
	+ \theta_{\kappa p\kappa}\, p^\kappa_t\right) d\kappa_t.
	\label{eq:32:dy_expanded}
\end{align}

Standard continuous-time growth accounting evaluates aggregate capital variation using a Divisia index of heterogeneous components \citep{Divisia1925}. Since $s_{\text{ME},t} + s_{\text{NRC},t} = 1$ and $d\kappa_t = dk^{\text{ME}}_t - dk^{\text{NRC}}_t$, aggregate capital growth decomposes as:
\begin{align}
	dk^{\text{cap}}_t &\approx s_{\text{ME},t}\, dk^{\text{ME}}_t
	+ s_{\text{NRC},t}\, dk^{\text{NRC}}_t
	\label{eq:32:dk_divisia} \\
	&= dk^{\text{NRC}}_t + s_{\text{ME},t}\, d\kappa_t.
	\label{eq:32:dk_simplified}
\end{align}
Dividing \eqref{eq:32:dy_expanded} by \eqref{eq:32:dk_simplified} yields the conventional total-derivative aggregate elasticity along the empirical growth path:
\begin{equation}
	\theta^{\text{Aggregate}}_{\text{total},t}
	= \frac{\theta_{\text{NRC}}
		+ \left(\theta_\kappa + \theta_{\omega\kappa}\, \omega_{t-1}
		+ \theta_{\kappa p\kappa}\, p^\kappa_t\right)
		\frac{d\kappa_t}{dk^{\text{NRC}}_t}}
	{1 + s_{\text{ME},t}\, \frac{d\kappa_t}{dk^{\text{NRC}}_t}}.
	\label{eq:32:theta_aggregate}
\end{equation}

Equation \eqref{eq:32:theta_aggregate} suffers from a severe theoretical flaw illuminated by the Cambridge Capital Controversy \citep{Robinson1953,Sraffa1960}. Heterogeneous physical assets cannot be aggregated into a single scalar $k^{\text{cap}}$ without value weights that depend on distribution and prices. Differentiating such a value aggregate and placing $dk^{\text{cap}}_t$ in the denominator re-introduces the circularity that the model with heterogeneous capital was designed to eliminate. Empirically, the denominator $1 + s_{\text{ME},t}(d\kappa_t / dk^{\text{NRC}}_t)$ vanishes whenever physical accumulation stalls or shifts composition, producing spurious zero-denominator singularities ($\theta \to \pm\infty$).

Alternatively, I construct the aggregate transformation elasticity as an asset weighted from the heterogeneous capital stock system. Rather than differentiating a value aggregate, the structural asset-weighted transformation elasticity weights identified heterogeneous elasticities by gross capital asset shares:
\begin{equation}
	\theta^{\text{Aggregate}}_t = \theta_{\text{NRC}}^{\text{str}} + s_{\text{ME},t} \cdot MP_{\kappa,t}^{\text{str}},
	\label{eq:32:theta_aggregate_structural}
\end{equation}
where $\theta_{\text{NRC}}^{\text{str}}$ is the base elasticity of structures, $s_{\text{ME},t} = K_{\text{ME},t} / K_{\text{cap},t}$ is the physical equipment share, and $MP_{\kappa,t}^{\text{str}} = \theta_{\kappa}^{\text{str}} + \theta_{\omega\kappa}\, (\omega_{t-1} - \bar{\omega}) + \theta_{\kappa p}\, (p^{\kappa}_t - \bar{p}^{\kappa})$ is the marginal capacity yield of mechanization. Equation \eqref{eq:32:theta_aggregate_structural} is immune to the Cambridge critique because the parameters are identified in a disaggregated system without prior scalar aggregation, and the weighting share $s_{\text{ME},t}$ is a physical asset ratio. This provides the equation for the trajectory to be implcitly identified from the co-integration relation estimated in empirically in Section~4.3.




\subsection{The Peripheral Extension: Balance-of-Payments
	Constraint and the Corner Solution}

In the capitalist periphery, the intensive margin of
accumulation is structurally tethered to the world market and
constrained by domestic social friction. Following
\citet{Prebisch1981} and \citet{Kaldor1959}, peripheral
capitalism is characterized by high consumption propensities
out of profits, structural inelasticities in agriculture, and
a chronic vulnerability to balance-of-payments crises
\citep{Ahumada2019}. This introduces a dual realization
bottleneck: the macroeconomic capacity to import and the
inelastic domestic supply of wage-goods. Section~2.5
established this constraint qualitatively as the peripheral
accumulation wedge. This subsection formalizes it as a strict
macroeconomic boundary constraint on the cost-minimizing
firm's choice of technique, integrating the structuralist
insights of \citet{Prebisch1981} into the heterogeneous
capital framework of Section~3.2.

\subsubsection{The Import-Biased Intensive Margin and the
	Balance of Payments Constraint}

In Section~3.2, the cost-minimizing firm chooses the
mechanization gap $\kappa_t$ to maximize the net viability
surplus:
\begin{equation}
	\min_{\kappa_t}\; c(\kappa_t)
	= \kappa_t \cdot MC(\omega_{t-1}, p^\kappa_t) - \Phi(\kappa_t)
	\label{eq:33:objective_periphery}
\end{equation}
where $\Phi(\kappa_t) = \lambda_0 + \lambda_1 \kappa_t
+ \lambda_2 \kappa_t^2$ is the regime-conditioned
mechanization frontier, and $MC(\omega_{t-1}, p^\kappa_t)
= \gamma_0 - \gamma_1 \omega_{t-1} + \gamma_2 p^\kappa_t$
is the marginal viability cost of the gap. In the core, this
optimization proceeds unconstrained, yielding the optimal
mechanization gap $\kappa_t^*$. In the periphery, however,
active mechanization is strictly import-biased: widening the
mechanization gap requires imported capital goods, specialized
machinery, and technologically non-substitutable spare parts \citep{Vernengo2006}.

\textbf{Assumption 3.} \textit{Import-biased mechanization.}
The physical import requirement per unit of output
($m^{req}_t$) is a linear function of the chosen technique:
\begin{equation}
	m^{req}_t = \nu \cdot \kappa_t
	\label{eq:33:import_req}
\end{equation}
where $\nu > 0$ captures the structural import-intensity of
the mechanization gap. Simultaneously, the physical capacity
to satisfy this requirement is constrained by the 
\textit{capacity to import} per unit of output
($m^{avail}_t = M^{avail}_t / Y_t$). Following
\citet{Prebisch1981}, peripheral export growth is governed by
the center's income elasticity of demand for primary goods
($\epsilon_{m,c} < 1$), structurally lagging behind domestic
import demand. This external asymmetry, compounded by high
consumption propensities out of capitalist profits
\citep{Kaldor1959}, severely squeezes the foreign exchange
available for productive accumulation. A necessary condition for the re-investment of the surplus. This imposes a strict
ceiling on the choice of technique: 
\begin{equation}
	\nu \cdot \kappa_t \le m^{avail}_t
	\implies \kappa_t \le \bar{\kappa}_t
	\equiv \frac{m^{avail}_t}{\nu}
	\label{eq:33:bop_ceiling}
\end{equation}

\subsubsection{The Macroeconomic Allocation of Foreign
	Exchange}

\textbf{Assumption 4.} \textit{Balance-of-payments
	constraint.} Total import demand cannot exceed export
capacity. The total supply of foreign exchange per unit of
output is driven by primary exports
($x_t = p_t^x X_t / Y_t$). Due to the center's low income
elasticity for primary goods, $x_t$ is structurally rigid.
The total demand for foreign exchange per unit of output
consists of three competing components:
\begin{enumerate}
	\item \textbf{Wage-goods imports} ($m^W_t$): driven by the
	wage share $\omega_{t-1}$ and the structural inelasticity
	of domestic agriculture \citep{Kaldor1959}.
	$m^W_t = m_w \omega_{t-1}$.
	\item \textbf{Conspicuous consumption imports} ($m^C_t$):
	driven by the profit share $\pi_t$ and the high propensity
	of the peripheral upper strata to imitate the consumption
	patterns of the center \citep{Prebisch1981, Kaldor1959}.
	$m^C_t = m_c c_k \pi_t$.
	\item \textbf{Productive capital imports} ($m^K_t$): driven
	by the mechanization gap. $m^K_t = \nu \kappa_t$.
\end{enumerate}

\textbf{Assumption 5.} \textit{Surplus leakage.} The
propensity to consume out of profits satisfies
$c_k > c_k^{\min}$, where $c_k^{\min}$ is the threshold
below which domestic savings suffice to finance capital
formation without external borrowing. Under the
developmentalist regime, the leakage channel is conspicuous
consumption imports by domestic elites
\citep{Prebisch1981, Kaldor1959}. Under the neoliberal
regime, the leakage channel is profit repatriation by
multinational corporations and capital flight
\citep{Ahumada2019}. In both cases, $c_k > c_k^{\min}$
ensures that the effective profit share available for
productive reinvestment is compressed below the nominal
level: $\pi_t^{eff} = \pi_t \cdot \phi_t < \pi_t$.

The balance of payments constraint requires:
\begin{equation}
	m_w \omega_{t-1} + m_c c_k \pi_t + \nu \kappa_t \le x_t
	\label{eq:33:bop_full}
\end{equation}
Using the distributive identity $\pi_t = 1 - \omega_{t-1}$,
I substitute and solve for the maximum feasible mechanization
gap:
\begin{align}
	m_w \omega_{t-1} + m_c c_k (1 - \omega_{t-1})
	+ \nu \kappa_t &\le x_t
	\label{eq:33:bop_sub} \\
	\nu \kappa_t &\le x_t - m_c c_k
	+ (m_c c_k - m_w) \omega_{t-1}
	\label{eq:33:bop_solved} \\
	\bar{\kappa}_t(\omega_{t-1})
	&= \frac{1}{\nu} \left[ x_t - m_c c_k
	+ (m_c c_k - m_w) \omega_{t-1} \right]
	\label{eq:33:kappa_bar}
\end{align}
This equation formalizes the structuralist critique: the base
capacity to import capital goods is determined by export
earnings minus the baseline luxury consumption of the
capitalist class ($x_t - m_c c_k$). The sign of
$(m_c c_k - m_w)$ determines whether a rising wage share
expands or contracts the feasible mechanization gap.

\subsubsection{The Constrained Cost-Surplus Maximization
	Problem}

The peripheral firm's tactical problem is therefore a cost-minimization problem constrained by the mechanization function and the external constraint, such that the minimization problem can be stated as a Kuhn–Tucker optimization problem: 
\begin{align}
	\min_{\kappa_t}\; c(\kappa_t)
	&=  \kappa_t \cdot MC(\omega_{t-1}, p^\kappa_t) - \Phi(\kappa_t)
	\label{eq:33:constrained_obj} \\
	\text{subject to:}\quad & \kappa_t \le \bar{\kappa}_t
	\label{eq:33:constraint}
\end{align}
This formulation distinct a regime shifting behavior of the choice of technique in the capitalist periphery given by a double constraint in the capitalists cost-minimizing behavior: i) Endogenous to domestic distributive conflict, ii) and to the dependency of foreign currency to fulfill import requirements at a given choice of technique. Therefore, if the external constraint becomes binding due lack of foreign currency, a regime shifting behavior of technological change takes place. 

\subsubsection{Endogenous Regime Shifts: The Corner Solution}

The firm's unconstrained optimal technique
$\kappa_t^*(\omega_{t-1})$ is strictly increasing in the wage
share ($\partial \kappa_t^* / \partial \omega_{t-1} > 0$) due
to the Okishio labor-saving bias established in
\eqref{eq:32:cs_omega}.

\textbf{Proposition 4.} \textit{Peripheral regime shifts.}
Two regimes are endogenous to the interaction between domestic
distributive conflict and the global market:
\begin{enumerate}
	\item \textbf{Regime A (Slack BoP):} When
	$\kappa_t^* \le \bar{\kappa}_t$, the firm solves the
	standard first-order condition, yielding $\kappa_t^*$. The
	labor-saving bias successfully extracts the Okishio
	distributive dividend.
	\item \textbf{Regime B (Peripheral realization trap):} When
	$\kappa_t^*(\omega_{t-1}) > \bar{\kappa}_t(\omega_{t-1})$,
	the constraint binds and the firm is forced into a corner
	solution:
	\begin{equation}
		\kappa_t = \bar{\kappa}_t
		\label{eq:33:corner}
	\end{equation}
\end{enumerate}
The corner solution arises through two intersecting structural
vulnerabilities. First, if $m_c c_k$ is large (the
conspicuous consumption drain), the baseline FX availability
is severely squeezed, lowering the intercept of
$\bar{\kappa}_t$. Second, if domestic agriculture is highly
inelastic ($m_w > m_c c_k$), a rising wage share
simultaneously increases the \textit{demand} for capital
imports (via $\kappa_t^*$) while \textit{contracting} the
supply of available FX (via $\bar{\kappa}_t$).
\textit{Proof.} Immediate from the Kuhn--Tucker conditions
applied to \eqref{eq:33:constrained_obj}--\eqref{eq:33:constraint}.
$\hfill\blacksquare$

\subsubsection{Effective Profit Compression and the Blocked
	Surplus}

\textbf{Definition 3.} \textit{FX-Conversion Ratio.} The
FX-conversion ratio $\phi_t$ is the ratio of the maximum
feasible technique to the optimal unconstrained technique:
\begin{equation}
	\phi_t = \min\left(1, \frac{\bar{\kappa}_t}{\kappa_t^*}\right)
	\label{eq:33:phi}
\end{equation}
When the external constraint binds ($\kappa_t^* > \bar{\kappa}_t$),
$\phi_t < 1$. This generates a direct compression of the
effective profit share available for intensive accumulation.

\textbf{Definition 4.} \textit{Blocked surplus.} Even if
domestic wage suppression yields a high nominal profit share
($\pi_t$), the physical inability to convert that surplus into
imported machinery squeezes the effective profit rate:
\begin{equation}
	\pi_t^{eff} = \pi_t \cdot \phi_t
	= (1 - \omega_{t-1}) \cdot \phi_t
	\label{eq:33:pi_eff}
\end{equation}
The difference constitutes the \textit{blocked surplus}:
\begin{equation}
	\text{Blocked Surplus} = \pi_t(1 - \phi_t)
	\label{eq:33:blocked}
\end{equation}
Because the accumulation of machinery is driven by the
\textit{effective} profit share, the
constraint decouples nominal profitability from capacity
formation. If capitalists squeeze wages to raise $\pi_t$, they
inadvertently increase conspicuous consumption ($m^C_t$),
which lowers $\bar{\kappa}_t$ and compresses $\phi_t$.

\subsubsection{The Regime-Shifting Aggregate
	Transformation Elasticity}

\textbf{Definition 5.} \textit{BoP binding ratio.} The BoP
binding ratio $\tau_t$ is the ratio of productive capital
imports to import capacity:
\begin{equation}
	\tau_t \equiv \frac{\nu \cdot I^{ME}_t}
	{\text{ToT}_t \times X_t}
	\label{eq:33:tau}
\end{equation}
where $I^{ME}_t$ is gross investment in machinery and equipment, and $\text{ToT}_t \times X_t$
proxies the import capacity (FX supply). Under Assumption~3,
productive capital imports $m^K_t = \nu \kappa_t$, since mechanization is import-biased in the periphery. The regime switch is governed by $\tau_t$:
\begin{itemize}
	\item $\tau_t < \bar{\tau}$: BoP constraint is
	\textbf{slack} (Regime A: unconstrained optimum).
	\item $\tau_t \ge \bar{\tau}$: BoP constraint is
	\textbf{binding} (Regime B: corner solution).
\end{itemize}
The threshold $\bar{\tau}$ is determined empirically by the
threshold regression in Section~4.

The peripheral extension generates a regime-shifting
aggregate transformation elasticity. Under \textbf{Regime~A}
(slack BoP), the constraint does not bind, and the
elasticity coincides with the unconstrained expression from
Section~3.2 for the aggregate elasticity:
\begin{equation}
	\theta^{A}_t = 
	= \frac{\theta_{NRC}
		+ \left(\theta_\kappa + \theta_{\omega\kappa}\omega_{t-1}
		+ \theta_{\kappa p\kappa} p^\kappa_t\right)
		\frac{d\kappa_t^*}{dk^{NRC}_t}}
	{1 + s_{ME,t}\frac{d\kappa_t^*}{dk^{NRC}_t}}
	\label{eq:33:theta_A}
\end{equation}

Under \textbf{Regime~B} (BoP binding), the surplus of capitalist enterprises becomes blocked, and the firm is forced
into a corner solution $\kappa_t = \bar{\kappa}_t$. The
marginal capacity payoff at the corner is strictly positive
because the firm would like to mechanize further but cannot:
\begin{equation}
	MP_\kappa\big|_{\kappa=\bar{\kappa}}
	= \Phi'(\bar{\kappa}_t) - MC(\omega_{t-1}, p^\kappa_t) > 0
	\label{eq:33:MP_corner}
\end{equation}
The Kuhn--Tucker multiplier on \eqref{eq:33:constraint} is
strictly positive, confirming that the constraint binds. The
aggregate transformation elasticity under the binding regime
becomes:
\begin{equation}
	\theta^{B}_t
	= \frac{\theta_{NRC}
		+ MP_\kappa\big|_{\kappa=\bar{\kappa}}
		\cdot \frac{d\bar{\kappa}_t}{dk^{NRC}_t}}
	{1 + s_{ME,t}\frac{d\bar{\kappa}_t}{dk^{NRC}_t}}
	\label{eq:33:theta_B}
\end{equation}
Since $\bar{\kappa}_t$ is determined by the BoP constraint
\eqref{eq:33:kappa_bar} rather than by the first-order
condition \eqref{eq:32:foc_kappa}, the growth path
$d\bar{\kappa}_t / dk^{NRC}_t$ differs from the unconstrained
path $d\kappa_t^* / dk^{NRC}_t$. The FX-conversion ratio
$\phi_t = \bar{\kappa}_t / \kappa_t^* < 1$ compresses the
intensive margin, and the elasticity under the binding regime
is strictly lower:
\begin{equation}
	\theta^{B}_t < \theta^{A}_t
	\label{eq:33:theta_inequality}
\end{equation}

\textbf{Proposition 5.} \textit{Regime-shifting elasticity.}
The peripheral aggregate transformation elasticity is a
regime-dependent parameter:
\begin{equation}
	\theta^{periphery}_t
	= \theta^{A}_t \cdot \mathbf{1}\{\kappa_t^* \le \bar{\kappa}_t\}
	+ \theta^{B}_t \cdot \mathbf{1}\{\kappa_t^* > \bar{\kappa}_t\}
	\label{eq:33:theta_regime}
\end{equation}
The elasticity shifts discretely when the BoP constraint
becomes binding. The regime switch is governed by the
interaction between the desired mechanization gap
$\kappa_t^*(\omega_{t-1}, p^\kappa_t)$ and the maximum
feasible gap $\bar{\kappa}_t(\omega_{t-1})$ determined by
export capacity, conspicuous consumption, and wage-goods
imports. An increase in the wage share $\omega_{t-1}$
simultaneously raises $\kappa_t^*$ (via capital-using 
labor-saving choice of technique) and can contract $\bar{\kappa}_t$ (via
the wage-goods import channel when $m_w > m_c c_k$), making
the regime switch more likely. This is what \citet{Prebisch1981} coined as the strangulation of economic development in the capitalist periphery: distributive pressures that would induce
productive mechanization in the center instead triggers the
BoP constraint, which in this framework explain the truncated or regime shifting behavior of the transformation elasticity.

\subsubsection{The Failure of the Viability Conditions for Technological Change and Structural Inflation}

\textbf{Proposition 6.} \textit{Blocked viability condition.}
At the boundary $\kappa_t = \bar{\kappa}_t$, the
implementation of the viability condition is blocked, i.e. technical change is not net-surplus increasing or conversely, inefficient because it fails to achieve a cost-minimizing choice of technique. The firm satisfies the viability criterion---the technique would reduce unit costs---but cannot deploy it because the required machinery cannot be imported due lack of foreign currency. Hence, the capitalist enterprise in the capitalist periphery suffers a persistent shortfall in its cost-surplus
relative to its unconstrained state.
\textit{Proof.} At $\kappa_t = \bar{\kappa}_t$, the
Kuhn--Tucker multiplier on \eqref{eq:33:constraint} is
strictly positive, so the first-order condition
$\partial c / \partial \kappa_t = 0$ does not hold. The
marginal viability yield exceeds the marginal cost at the
boundary, but the firm cannot expand further.
$\hfill\blacksquare$


\textbf{Corollary 1.} \textit{Structural Inflation} Unable to extract the distributive dividend via mechanization, capitalists attempt to recover the blocked surplus through the only remaining channel: price markups. When the wage-goods bottleneck forces nominal wage adjustments and the external BoP
bottleneck blocks defensive mechanization, any further increase
in the wage share directly crushes profitability. The resulting
defensive pricing behavior ignites structural inflation
dynamic \citep{Kalecki1955, Sunkel1958, Prebisch1981}. This is the
peripheral realization trap: nominal profitability coexists
with truncated mechanization and structurally stalled built-environments that constitute the \textit{``envelope''} of the labor process,
permanently destabilizing the long-run relation between investments in plant 
and machinery.


\subsection{Analytical and Theoretical Synthesis}

The three models developed in this section formalize a single
theoretical claim: that decentralized, individually
cost-minimizing decisions generate aggregate
disproportionality under the unbalanced growth closure.
Table~\ref{tab:coherence} summarizes the articulation
between the Okishio and Vidal frameworks: 

\begin{table}[H]
	\centering
	\footnotesize
	\caption{Market Anarchy and the Internally Divided Firm Articulation Across Three Models}
	\label{tab:coherence}
	\begin{tabular}{p{2.2cm}p{3.8cm}p{3.8cm}p{4.2cm}}
		\toprule
		Logic & Decentralized Accumulation &
		Internally Divided Firm & Analytical Insight \\
		\midrule
		\midrule
		Micro logic & Viability condition (cost
		reduction at prevailing prices) & Internal
		contradiction (control vs.\ cooperation) &
		Optimization problem = cost-minimizing pole;
		constraints = cooperation limits \\
		Macro logic & Decentralized decisions $\to$
		disproportionality & Divided firm $\to$
		contested operating conditions & $\theta \neq 1$
		captures gap between accumulation and capacity \\
		Disequilibrium & Unbalanced growth (supply
		$\neq$ demand) & ``Mediocre sufficiency''
		\citep{Vidal2022} (no coherent optimum) & $\mu$
		demonstrates gap; $\theta$ captures structural
		inefficiency \\
		Peripheral specificity & BoP constraint blocks
		\textit{implementation} of viability & External
		constraint compounds internal division & Corner
		solution: $\kappa = \bar{\kappa}$ when
		constraint binds \\
		\bottomrule
	\end{tabular}
\end{table}

This analytical framework does not assert equilibrium. It aims 
demonstrates the opposite: capitalist accumulation,
left to its own decentralized logic, does not generate
proportional reproduction. The rate of capacity utilization 
$\mu$ is the exhaust valve of this
disproportionality: it shows that the normal operating
conditions of capitalist firms, determined by their
cost-minimizing mechanization choices, do not meet
aggregate demand conditions. The transformation elasticity
$\theta$ captures this structural gap between accumulation
and realized capacity. Whether $\theta$ exceeds or falls
below unity is an empirical question to be tackled below. In the center, this gap is cyclical---mediated by institutional configurations that
can be exhausted triggering processes of capitalist re-structuring. In the periphery, this gap is permanent---locked in by the structural ceiling of
external subordination \citep{Kaldor1959,Prebisch1981}.

Three models have now specified what $\theta$ depends on, and none of
them has said what it equals. The dependence is on objects that no
national accounting system reports directly: a capacity ceiling that is
never observed, a mechanization gap between two capital stocks that are
valued by convention, and, for the periphery, an external state that has
to be inferred rather than read off a date. Whether $\theta$ can be
recovered at all therefore turns on whether those objects can be
reconstructed from the historical record without importing the
transformation relation into their construction.